% CS301 Week 11 -- Trees: terminology; binary trees and their properties;
% tree traversals; expression trees; binary search trees (search, insert,
% delete, min/max); AVL trees and rotations; Unit-III review.
% Difficulty: intro (course-wide default per knowledge-base-CS301.md).
% Written in plain language with small hand-traced examples.
% Page-grounded in Wirth2004 Sec.4.4 (Tree Structures, p.126), Sec.4.4.2
% (Basic Operations, p.132), Sec.4.4.3 (Tree Search and Insertion, p.135),
% Sec.4.4.4 (Tree Deletion, p.139), Sec.4.5 (Balanced Trees, p.142);
% Sedgewick2011 Sec.3.2 (Binary Search Trees, p.396), Sec.3.3 (Balanced
% Search Trees, p.424, for the balancing motivation; its balanced tree is
% red-black, not AVL); Karumanchi2017 Ch.6 (Trees, pp.224-368 PDF).
% Horowitz & Sahni Ch.5 and Aho-Hopcroft-Ullman Ch.3, 5 are cited at chapter
% level only (not indexed) -- no page numbers are asserted for them.
% Handoff: Week 10 closed on "hashing finds a key in one probe but cannot
% answer the next key or all keys between a and b. Week 11 builds trees --
% binary trees, BSTs and AVL rotations -- that keep order and logarithmic
% search." This deck pays that off with the calculator's symbol table: the
% same 2,000 identifiers, now printable in sorted order and searchable in
% O(log n) even while inserts continue.
% Sets up Week 12 (graphs): a tree is a graph with no cycles and one root.

\documentclass[aspectratio=169]{beamer}
\input{header}
\coursecode{CS301}

% Shared tree-drawing styles. Explicit sizes keep the P7 clearance audit easy:
% every node has a declared minimum size, and each diagram has a coordinate map.
\tikzset{
  tnode/.style   = {circle, draw=primary-blue, fill=light-bg,
                    minimum size=0.80cm, inner sep=0pt, font=\small},
  tnodeg/.style  = {circle, draw=primary-gold, line width=1.1pt,
                    fill=highlight-yellow!25, minimum size=0.80cm,
                    inner sep=0pt, font=\small},
  tnodesm/.style = {circle, draw=primary-blue, fill=light-bg,
                    minimum size=0.55cm, inner sep=0pt, font=\scriptsize},
  tnodesmg/.style= {circle, draw=primary-gold, line width=1pt,
                    fill=highlight-yellow!25, minimum size=0.55cm,
                    inner sep=0pt, font=\scriptsize},
  tedge/.style   = {draw=neutral, thick}
}

\title{Keeping Order}
\subtitle{Week 11 --- Trees, Binary Search Trees \& AVL}
\author[Dr. Auqib Hamid Lone]{\textbf{Dr. Auqib Hamid Lone}}
\institute[GCET Ganderbal]{PCC CS-301: Data Structures\\GCET, Ganderbal}
\date[Week 11]{\small Week 11}

\begin{document}

\frame{\titlepage}

% =============================================================================
% ACT 0 --- MOTIVATION + HANDOFF
% =============================================================================
\section{Motivation and the Handoff}

\begin{frame}{Today's Five Questions}
  \begin{enumerate}
    \item What is a tree, in words a first-year student can draw?
    \item What are the four standard ways to visit every node, and what is
      each one for?
    \item How does a binary search tree keep keys ordered while they are
      still being inserted and deleted?
    \item Why can a plain BST slow down to a linked list --- and what does an
      AVL tree do about it?
    \item When is a tree the right choice, and when is a hash table still
      better?
  \end{enumerate}
  \vspace{0.15cm}
  \begin{exampleblock}{Plan}
    Learn the vocabulary on one small tree. Walk it four ways. Turn an
    expression into a tree. Then add one ordering rule to get a BST, and one
    balancing rule to keep it fast.
  \end{exampleblock}
\end{frame}

\begin{frame}{Handoff: What Week 10 Left Open}
  \begin{itemize}
    \item Week 10 found a record in about one probe with a hash table.
      It closed with a promise: hashing cannot answer ``the next key'' or
      ``all keys between $a$ and $b$''; Week 11 builds trees that keep
      \emph{order} and \emph{logarithmic} search.
    \item The bill Week 10 left unpaid:
      \begin{itemize}
        \item a sorted array searches in $O(\log n)$ but inserts in $O(n)$;
        \item a hash table inserts in $O(1)$ expected but keeps no order;
        \item a linked list keeps insertion order but searches in $O(n)$.
      \end{itemize}
    \item New words today: \key{root}, \key{leaf}, \key{height},
      \key{traversal}, \key{binary search tree (BST)}, and
      \key{balance factor}.
  \end{itemize}
  \begin{exampleblock}{This week's job}
    Build a shape that keeps keys sorted \emph{and} keeps search short while
    the data changes. A tree gives us both --- if we keep its height small.
  \end{exampleblock}
\end{frame}

\begin{frame}{What We Want From a Dictionary}
  \begin{center}
  {\footnotesize
  \begin{tabular}{l c c c}
    \toprule
    \textbf{structure} & \textbf{search} & \textbf{insert} & \textbf{ordered?} \\
    \midrule
    sorted array (Weeks 8--9) & $O(\log n)$ & $O(n)$ & yes \\
    hash table (Week 10)      & $O(1)$ expected & $O(1)$ expected & no \\
    linked list (Week 6)      & $O(n)$ & $O(1)$ at head & no \\
    \midrule
    balanced BST (this week)  & $O(\log n)$ & $O(\log n)$ & yes \\
    \bottomrule
  \end{tabular}}
  \end{center}
  \begin{itemize}
    \item We want the last row: fast search, fast insert, and order.
    \item A tree is the first structure in this course that can give all
      three at once. The price is that we must keep its \emph{shape} under
      control.
  \end{itemize}
  \muted{Running application (Weeks 3--11): the calculator's symbol table.
    Week 10 could find an identifier fast but could not print the identifiers
    in alphabetical order. This week's tree can do both.}
\end{frame}

\transitionslide{A list is a line. A tree is a shape.}

% =============================================================================
% ACT 1 --- TREE VOCABULARY
% =============================================================================
\section{Tree Vocabulary}

\sectiondivider{Section 1}{Nodes, Edges, Depth, Height}

\begin{frame}{One Picture, Seven Words}
  \begin{center}
    % Coordinate map: root (0,3); depth-1 nodes (+-1.7,1.8);
    % depth-2 leaves (+-2.5 or +-0.9, 0.6). Node size 0.8 cm.
    \begin{tikzpicture}[every node/.style={font=\small}]
      \node[tnode] (n50) at (0,3)   {50};
      \node[tnode] (n30) at (-1.7,1.8) {30};
      \node[tnode] (n70) at (1.7,1.8)  {70};
      \node[tnode] (n20) at (-2.5,0.6) {20};
      \node[tnode] (n40) at (-0.9,0.6) {40};
      \node[tnode] (n60) at (0.9,0.6)  {60};
      \node[tnode] (n80) at (2.5,0.6)  {80};
      \draw[tedge] (n50) -- (n30);
      \draw[tedge] (n50) -- (n70);
      \draw[tedge] (n30) -- (n20);
      \draw[tedge] (n30) -- (n40);
      \draw[tedge] (n70) -- (n60);
      \draw[tedge] (n70) -- (n80);
      \foreach \y/\t in {3/0, 1.8/1, 0.6/2}
        \node[text=neutral, anchor=east, font=\scriptsize] at (-3.7,\y)
          {depth \t};
    \end{tikzpicture}
  \end{center}
  \vspace{-0.1cm}
  \footnotesize
  \begin{itemize}
    \item \key{Root}: 50 --- the top node, with no parent.
    \item \key{Leaf}: 20, 40, 60, 80 --- nodes with no children.
    \item \key{Parent} of 20 is 30; \key{child} of 30 is 20;
      \key{siblings} are 20 and 40.
    \item A \key{subtree} is any node plus everything below it: 30, 20, 40
      is the left subtree of 50.
  \end{itemize}
\end{frame}

\begin{frame}{Depth and Height Are Counted in Edges}
  \vspace{-0.15cm}
  \begin{itemize}
    \item \key{Depth} of a node = number of edges from the root down to it.
      The root has depth $0$.
    \item \key{Height} of a node = number of edges on the longest path from
      it down to a leaf. A leaf has height $0$.
    \item \key{Height of the tree} = height of its root.
  \end{itemize}
  \begin{center}
  {\footnotesize
  \begin{tabular}{l c c}
    \toprule
    \textbf{node} & \textbf{depth} & \textbf{height} \\
    \midrule
    50 & 0 & 2 \\
    30 & 1 & 1 \\
    70 & 1 & 1 \\
    20, 40, 60, 80 & 2 & 0 \\
    \bottomrule
  \end{tabular}}
  \end{center}
  \begin{exampleblock}{Convention used in this course}
    A single node has height $0$; the empty tree has height $-1$. State the
    convention when you write it down --- books differ, and an exam answer
    must say which one it uses.
  \end{exampleblock}
  \muted{Tree definitions follow \cite[pp.~126--131]
    {Wirth2004_algorithms_data_structures}.}
\end{frame}

\begin{frame}{A Binary Tree Has at Most Two Children}
  \begin{block}{Definition --- binary tree}
    A \key{binary tree} is a tree in which every node has at most two
    children, called its \key{left child} and \key{right child}.
  \end{block}
  \begin{itemize}
    \item The two sides are \emph{not} interchangeable. Swapping a left and
      right child changes the tree, even when the drawing looks similar.
    \item Our seven-node tree is a binary tree. A node may have $0$, $1$ or
      $2$ children --- never three.
    \item In C:
      \texttt{struct node \{ int data; struct node *left, *right; \};}
  \end{itemize}
  \muted{The pointer picture is the Week 6 linked-list node with two links
    instead of one: \texttt{left} and \texttt{right}.}
\end{frame}

\begin{frame}{Three Shapes Worth Naming}
  \begin{center}
    % Left panel: perfect tree (7 nodes, height 2). Right panel: the same
    % tree missing node 7, which is complete but not perfect.
    \begin{tikzpicture}[every node/.style={font=\scriptsize}]
      \foreach \x/\y/\v in {1/2/4, 0.4/1/2, 1.6/1/6,
                            0.1/0/1, 0.7/0/3, 1.3/0/5, 1.9/0/7}
        \node[tnodesm] at (\x,\y) {\v};
      \draw[tedge] (1,2) -- (0.4,1);
      \draw[tedge] (1,2) -- (1.6,1);
      \draw[tedge] (0.4,1) -- (0.1,0);
      \draw[tedge] (0.4,1) -- (0.7,0);
      \draw[tedge] (1.6,1) -- (1.3,0);
      \draw[tedge] (1.6,1) -- (1.9,0);
      \node[text=neutral, font=\tiny] at (1,-0.75)
        {perfect: 7 nodes, leaves at depth 2};

      \foreach \x/\y/\v in {5.5/2/4, 4.9/1/2, 6.1/1/6,
                            4.6/0/1, 5.2/0/3, 5.8/0/5}
        \node[tnodesm] at (\x,\y) {\v};
      \draw[tedge] (5.5,2) -- (4.9,1);
      \draw[tedge] (5.5,2) -- (6.1,1);
      \draw[tedge] (4.9,1) -- (4.6,0);
      \draw[tedge] (4.9,1) -- (5.2,0);
      \draw[tedge] (6.1,1) -- (5.8,0);
      \node[text=neutral, font=\tiny] at (5.5,-0.75)
        {complete: last level fills left to right};
    \end{tikzpicture}
  \end{center}
  \vspace{-0.2cm}
  \footnotesize
  \begin{itemize}
    \item \key{Full}: every node has $0$ or $2$ children. The seven-node
      tree above is full.
    \item \key{Complete}: every level is full except the last, and the last
      level is filled left to right.
    \item \key{Perfect}: full and all leaves are at the same depth. The
      left-hand tree is perfect.
  \end{itemize}
\end{frame}

\begin{frame}{Counting Nodes: Four Useful Facts}
  \begin{itemize}
    \item Level $d$ holds at most $2^d$ nodes (root is level $0$):
      $1, 2, 4, 8, \dots$
    \item A tree of height $h$ holds at most
      $1 + 2 + 4 + \cdots + 2^h = 2^{h+1} - 1$ nodes.
    \item Turn that around: a tree with $n$ nodes has height at least
      $\lfloor \log_2 n \rfloor$. A perfect tree with $n$ nodes has height
      exactly $\log_2(n+1) - 1$.
    \item In any binary tree with $n$ nodes there are exactly $n+1$ empty
      child slots (\texttt{NULL} pointers). A full binary tree with $L$
      leaves has $L-1$ internal nodes.
  \end{itemize}
  \begin{exampleblock}{Why the logarithm appears}
    Each level can double the number of nodes, so height grows like
    $\log_2 n$. A tree with $1{,}023$ nodes can be only $9$ edges tall.
    That is the whole reason trees can search quickly --- if we keep them
    short.
  \end{exampleblock}
\end{frame}

\begin{frame}{Same Keys, Different Shape, Different Speed}
  \begin{center}
  {\footnotesize
  \begin{tabular}{l c c}
    \toprule
    \textbf{shape of the tree} & \textbf{height} & \textbf{comparisons to find 50} \\
    \midrule
    a chain (inserted in sorted order) & 4 & 5 \\
    a balanced shape & 2 & 3 \\
    \bottomrule
  \end{tabular}}
  \end{center}
  \begin{itemize}
    \item The keys are the same $\{10,20,30,40,50\}$. Only the shape changed.
    \item Search cost is the number of nodes on the root-to-target path,
      which is at most height $+1$.
    \item So the real question for this week is not ``what is a tree?'' but
      \key{how do we keep the tree short?}
  \end{itemize}
\end{frame}

\transitionslide{Every node can be visited in more than one order.}

% =============================================================================
% ACT 2 --- TRAVERSALS
% =============================================================================
\section{Tree Traversals}

\sectiondivider{Section 2}{Four Ways to Visit Every Node}

\begin{frame}{The Tree We Will Walk}
  \begin{center}
    \begin{tikzpicture}[every node/.style={font=\small}]
      \node[tnode] (n50) at (0,2.8)   {50};
      \node[tnode] (n30) at (-1.7,1.7) {30};
      \node[tnode] (n70) at (1.7,1.7)  {70};
      \node[tnode] (n20) at (-2.5,0.6) {20};
      \node[tnode] (n40) at (-0.9,0.6) {40};
      \node[tnode] (n60) at (0.9,0.6)  {60};
      \node[tnode] (n80) at (2.5,0.6)  {80};
      \draw[tedge] (n50) -- (n30);
      \draw[tedge] (n50) -- (n70);
      \draw[tedge] (n30) -- (n20);
      \draw[tedge] (n30) -- (n40);
      \draw[tedge] (n70) -- (n60);
      \draw[tedge] (n70) -- (n80);
    \end{tikzpicture}
  \end{center}
  \begin{itemize}
    \item A \key{traversal} visits every node exactly once and prints its
      key in some order.
    \item The order is not arbitrary: it is chosen by \emph{when} we print
      the node relative to its two subtrees.
  \end{itemize}
  \muted{All four traversals cost $O(n)$ time: each node is handled once.
    They differ only in the order of the output. Definitions follow
    \cite[pp.~132--135]{Wirth2004_algorithms_data_structures} and
    \cite[pp.~396--423]{Sedgewick2011_algorithms}.}
\end{frame}

\begin{frame}{Inorder: Left, Node, Right}
  \begin{block}{Rule}
    Traverse the left subtree, then print the node, then traverse the right
    subtree.
  \end{block}
  \begin{itemize}
    \item Start at 50: go left to 30, go left to 20. 20 has no left child,
      so print \textbf{20}; no right child, return.
    \item Back at 30: print \textbf{30}; go right to 40: print \textbf{40}.
    \item Back at 50: print \textbf{50}; then the same walk on the right
      subtree gives \textbf{60, 70, 80}.
  \end{itemize}
  \begin{exampleblock}{Output}
    $20,\ 30,\ 40,\ 50,\ 60,\ 70,\ 80$ --- sorted!
  \end{exampleblock}
  \muted{That is not a coincidence. In a BST, inorder visits keys from
    smallest to largest; we will use this as a check in Section 4.}
\end{frame}

\begin{frame}{Preorder: Node, Left, Right}
  \begin{block}{Rule}
    Print the node first, then traverse the left subtree, then the right
    subtree.
  \end{block}
  \begin{itemize}
    \item Print \textbf{50}; go left: print \textbf{30}; print \textbf{20};
      return and print \textbf{40}.
    \item Return to 50; go right: print \textbf{70}, then \textbf{60}, then
      \textbf{80}.
  \end{itemize}
  \begin{exampleblock}{Output}
    $50,\ 30,\ 20,\ 40,\ 70,\ 60,\ 80$
  \end{exampleblock}
  \begin{itemize}
    \item Use it to \emph{copy} a tree or write it to a file: the parent is
      always printed before its children, so the same shape can be rebuilt.
    \item It is also the prefix form of an expression (Section 3).
  \end{itemize}
\end{frame}

\begin{frame}{Postorder: Left, Right, Node}
  \begin{block}{Rule}
    Traverse both subtrees first, then print the node.
  \end{block}
  \begin{itemize}
    \item Left subtree of 50: visit 20, then 40, then \textbf{30}.
    \item Right subtree: visit 60, then 80, then \textbf{70}.
    \item Finally print the root \textbf{50}.
  \end{itemize}
  \begin{exampleblock}{Output}
    $20,\ 40,\ 30,\ 60,\ 80,\ 70,\ 50$
  \end{exampleblock}
  \begin{itemize}
    \item Use it to \emph{free} a tree: children are freed before their
      parent, so no live pointer is left dangling.
    \item It is also how an expression tree is evaluated (Section 3).
  \end{itemize}
\end{frame}

\begin{frame}{Level Order: Top to Bottom, Left to Right}
  \begin{block}{Rule}
    Visit the root, then all nodes at depth $1$, then all at depth $2$, and
    so on; within a level, go left to right.
  \end{block}
  \begin{itemize}
    \item Depth $0$: \textbf{50}. Depth $1$: \textbf{30, 70}.
      Depth $2$: \textbf{20, 40, 60, 80}.
    \item This is \emph{not} recursion on the two children. Use the queue
      from Week 5: enqueue the root; while the queue is non-empty, dequeue a
      node, print it, and enqueue its children.
  \end{itemize}
  \begin{exampleblock}{Output}
    $50,\ 30,\ 70,\ 20,\ 40,\ 60,\ 80$
  \end{exampleblock}
  \muted{Level order is breadth-first search on a tree --- the same idea
    returns for graphs in Week 12.}
\end{frame}

\begin{frame}{The Same Tree, Four Orders}
  \begin{center}
  {\footnotesize
  \begin{tabular}{l l}
    \toprule
    \textbf{traversal} & \textbf{output for the seven-node tree} \\
    \midrule
    inorder   & 20, 30, 40, 50, 60, 70, 80 \\
    preorder  & 50, 30, 20, 40, 70, 60, 80 \\
    postorder & 20, 40, 30, 60, 80, 70, 50 \\
    level order & 50, 30, 70, 20, 40, 60, 80 \\
    \bottomrule
  \end{tabular}}
  \end{center}
  \begin{itemize}
    \item Inorder is the only one that reads this particular tree in sorted
      order --- because this tree happens to be a BST.
    \item A quick memory hook: \key{in}order puts the node \emph{in} the
      middle; \key{pre}order puts it before; \key{post}order puts it after.
  \end{itemize}
\end{frame}

\begin{frame}[fragile]{Traversals in C}
  \scriptsize
\begin{verbatim}
void inorder(struct node *t) {
    if (t == NULL) return;
    inorder(t->left);
    printf("%d ", t->data);
    inorder(t->right);
}

void preorder(struct node *t) {
    if (t == NULL) return;
    printf("%d ", t->data);
    preorder(t->left);
    preorder(t->right);
}

void postorder(struct node *t) {
    if (t == NULL) return;
    postorder(t->left);
    postorder(t->right);
    printf("%d ", t->data);
}
\end{verbatim}
  \vspace{-0.15cm}
  \footnotesize
  \muted{The three functions differ only in where \texttt{printf} sits;
    each visits every node once, so each is $O(n)$.}
\end{frame}

\begin{frame}{Check Yourself: Walk a New Tree}
  \begin{center}
    \begin{tikzpicture}[every node/.style={font=\scriptsize}]
      \node[tnodesm] (a) at (0,2.6) {8};
      \node[tnodesm] (b) at (-1.4,1.5) {3};
      \node[tnodesm] (c) at (1.4,1.5) {10};
      \node[tnodesm] (d) at (-2.1,0.4) {1};
      \node[tnodesm] (e) at (-0.7,0.4) {6};
      \node[tnodesm] (f) at (1.4,0.4) {14};
      \node[tnodesm] (g) at (-1.1,-0.7) {4};
      \node[tnodesm] (h) at (-0.3,-0.7) {7};
      \node[tnodesm] (i) at (0.7,0.4) {13};
      \draw[tedge] (a) -- (b);  \draw[tedge] (a) -- (c);
      \draw[tedge] (b) -- (d);  \draw[tedge] (b) -- (e);
      \draw[tedge] (c) -- (f);  \draw[tedge] (e) -- (g);
      \draw[tedge] (e) -- (h);  \draw[tedge] (f) -- (i);
    \end{tikzpicture}
  \end{center}
  \vspace{-0.25cm}
  \small
  \begin{itemize}
    \item Write the three outputs yourself, then check the answers below.
    \item Inorder: 1, 3, 4, 6, 7, 8, 10, 13, 14.
    \item Preorder: 8, 3, 1, 6, 4, 7, 10, 14, 13.
    \item Postorder: 1, 4, 7, 6, 3, 13, 14, 10, 8.
  \end{itemize}
  \muted{If your inorder is not sorted, this tree is not a BST --- check
    the ordering rule in Section 4.}
\end{frame}

\transitionslide{An expression is a tree in disguise.}

% =============================================================================
% ACT 3 --- EXPRESSION TREES
% =============================================================================
\section{Expression Trees}

\sectiondivider{Section 3}{From Postfix to a Tree and Back}

\begin{frame}{Building an Expression Tree from Postfix}
  \begin{itemize}
    \item Week 4 converted \texttt{(a+b)*c} to postfix \texttt{a b + c *}.
      Now read that postfix left to right with a stack of \emph{trees}:
    \item \textbf{Operand} (a number or variable): make a one-node tree and
      push it.
    \item \textbf{Operator}: pop the \emph{right} tree, then the \emph{left}
      tree; make a new node holding the operator, with those two trees as
      children; push the new tree.
    \item At the end, the stack holds exactly one tree: the expression.
  \end{itemize}
  \begin{exampleblock}{Trace of \texttt{a b + c *}}
    push \texttt{a}; push \texttt{b}; \texttt{+} pops \texttt{b} (right) and
    \texttt{a} (left) $\to$ node \texttt{+}. push \texttt{c}; \texttt{*}
    pops \texttt{c} (right) and \texttt{+} (left) $\to$ root \texttt{*}. \\
    \good{Pop right before left} --- reversing this swap gives a wrong tree.
  \end{exampleblock}
  \muted{Expression trees: \cite[Ch.~6, pp.~224--368 (PDF)]
    {Karumanchi2017_data_structures_algorithms_made_easy}.}
\end{frame}

\begin{frame}{The Tree for \texttt{(a+b)*c}}
  \begin{center}
    \begin{tikzpicture}[every node/.style={font=\small}]
      \node[tnode] (star) at (0,2.4) {\texttt{*}};
      \node[tnode] (plus) at (-1.2,1.2) {\texttt{+}};
      \node[tnode] (c)    at (1.2,1.2)  {\texttt{c}};
      \node[tnode] (a)    at (-2.0,0.1) {\texttt{a}};
      \node[tnode] (b)    at (-0.4,0.1) {\texttt{b}};
      \draw[tedge] (star) -- (plus);
      \draw[tedge] (star) -- (c);
      \draw[tedge] (plus) -- (a);
      \draw[tedge] (plus) -- (b);
    \end{tikzpicture}
  \end{center}
  \begin{itemize}
    \item Operators are \emph{internal} nodes; operands are \emph{leaves}.
    \item The root is the operator applied last. Its left subtree computes
      \texttt{a+b}; its right subtree is \texttt{c}.
    \item Read the leaves left to right: \texttt{a b + c *} --- the postfix
      form comes back for free.
  \end{itemize}
\end{frame}

\begin{frame}{Evaluating the Tree: Postorder}
  \begin{itemize}
    \item To evaluate, evaluate the left subtree, evaluate the right
      subtree, then apply the operator. That is exactly postorder.
    \item Trace with $a=2$, $b=3$, $c=4$:
      \begin{itemize}
        \item leaf \texttt{a} $\to 2$; leaf \texttt{b} $\to 3$;
        \item node \texttt{+} $\to 2+3=5$;
        \item leaf \texttt{c} $\to 4$;
        \item root \texttt{*} $\to 5\times 4=20$.
      \end{itemize}
  \end{itemize}
  \begin{exampleblock}{Output}
    $20$ --- and no parentheses were needed anywhere in the tree.
  \end{exampleblock}
  \muted{The tree stores the grouping that parentheses would otherwise have
    to spell out. That is why compilers build trees.}
\end{frame}

\begin{frame}{Inorder Gives Infix Back --- With a Warning}
  \begin{itemize}
    \item Inorder on the same tree prints
      $\texttt{a} + \texttt{b} \times \texttt{c}$.
    \item That is \emph{not} the original expression! Without parentheses,
      the usual precedence rule reads it as \texttt{a + (b*c)}.
    \item To recover the meaning, put parentheses around every operator's
      result as you return:
      $(\texttt{a}+\texttt{b})$, then $((\texttt{a}+\texttt{b})\times
      \texttt{c})$.
    \item The three walks match the three notations from Week 4:
  \end{itemize}
  \begin{center}
  {\footnotesize
  \begin{tabular}{l l}
    \toprule
    \textbf{walk} & \textbf{notation} \\
    \midrule
    inorder   & infix (parentheses needed) \\
    preorder  & prefix \\
    postorder & postfix (and evaluation order) \\
    \bottomrule
  \end{tabular}}
  \end{center}
\end{frame}

\transitionslide{Now add one rule that makes search fast.}

% =============================================================================
% ACT 4 --- BINARY SEARCH TREES
% =============================================================================
\section{Binary Search Trees}

\sectiondivider{Section 4}{The BST Rule: Left Is Smaller}

\begin{frame}{The BST Rule}
  \begin{block}{Definition --- binary search tree}
    For \emph{every} node $x$: every key in the left subtree of $x$ is
    smaller than $x$'s key, and every key in the right subtree is larger.
  \end{block}
  \begin{itemize}
    \item The rule must hold at \emph{every} node, not just between a parent
      and its immediate child.
    \item This is what makes search possible: at each node, one comparison
      throws away an entire subtree.
    \item Duplicates: pick a policy and state it. Simplest is to reject a
      duplicate; another is to keep a count in the node.
  \end{itemize}
  \muted{A binary tree plus this one rule $=$ a BST. The shape alone does not
    make a BST. Search and insertion: \cite[pp.~135--139]
    {Wirth2004_algorithms_data_structures}; \cite[pp.~396--423]
    {Sedgewick2011_algorithms}; chapter-level treatment in
    \cite[Ch.~5]{HorowitzSahni2008_fundamentals_data_structures}.}
\end{frame}

\begin{frame}{Is This a BST?}
  \begin{center}
    % Left: valid. Right: node 6 sits in the left subtree of 5 but is > 5.
    \begin{tikzpicture}[every node/.style={font=\scriptsize}]
      \node[tnodesm] (a) at (0,2.2) {5};
      \node[tnodesm] (b) at (-0.9,1.2) {3};
      \node[tnodesm] (c) at (0.9,1.2) {8};
      \node[tnodesm] (d) at (-1.4,0.2) {1};
      \node[tnodesm] (e) at (-0.4,0.2) {4};
      \node[tnodesm] (f) at (0.4,0.2) {7};
      \node[tnodesm] (g) at (1.4,0.2) {9};
      \draw[tedge] (a)--(b); \draw[tedge] (a)--(c);
      \draw[tedge] (b)--(d); \draw[tedge] (b)--(e);
      \draw[tedge] (c)--(f); \draw[tedge] (c)--(g);
      \node[text=positive, font=\scriptsize] at (0,-0.55) {\good{yes}};

      \node[tnodesm] (A) at (4.5,2.2) {5};
      \node[tnodesm] (B) at (3.6,1.2) {3};
      \node[tnodesm] (C) at (5.4,1.2) {8};
      \node[tnodesm] (D) at (3.1,0.2) {1};
      \node[tnodesmg] (E) at (4.1,0.2) {6};
      \node[tnodesm] (F) at (4.9,0.2) {7};
      \node[tnodesm] (G) at (5.9,0.2) {9};
      \draw[tedge] (A)--(B); \draw[tedge] (A)--(C);
      \draw[tedge] (B)--(D); \draw[tedge] (B)--(E);
      \draw[tedge] (C)--(F); \draw[tedge] (C)--(G);
      \node[text=negative, font=\scriptsize] at (4.5,-0.55)
        {\bad{no: 6 $>$ 5 but lies left of 5}};
    \end{tikzpicture}
  \end{center}
  \vspace{-0.2cm}
  \begin{itemize}
    \item In the right-hand tree, 6 is a child of 3, and $6>3$ looks fine.
      The violation is against the \emph{root} 5: 6 is in 5's left subtree.
    \item Moral: checking only parent--child pairs is not enough.
  \end{itemize}
\end{frame}

\begin{frame}{Search: Follow the Comparisons}
  \begin{center}
    \begin{tikzpicture}[every node/.style={font=\small}]
      \node[tnode] (n50) at (0,2.8)   {50};
      \node[tnode] (n30) at (-1.7,1.7) {30};
      \node[tnode] (n70) at (1.7,1.7)  {70};
      \node[tnode] (n20) at (-2.5,0.6) {20};
      \node[tnodeg] (n40) at (-0.9,0.6) {40};
      \node[tnode] (n60) at (0.9,0.6)  {60};
      \node[tnode] (n80) at (2.5,0.6)  {80};
      \draw[tedge] (n50) -- (n30);
      \draw[tedge] (n50) -- (n70);
      \draw[tedge] (n30) -- (n20);
      \draw[tedge] (n30) -- (n40);
      \draw[tedge] (n70) -- (n60);
      \draw[tedge] (n70) -- (n80);
      \draw[-{Latex[length=2mm]}, primary-gold, thick]
        (n50) to[bend left=12] (n30);
      \draw[-{Latex[length=2mm]}, primary-gold, thick]
        (n30) to[bend left=12] (n40);
    \end{tikzpicture}
  \end{center}
  \vspace{-0.15cm}
  \begin{itemize}
    \item Search for 40: $40<50$ go left; $40>30$ go right; equal --- found.
      Three comparisons, not seven.
    \item At each node: equal $\to$ found; smaller $\to$ left; larger
      $\to$ right; \texttt{NULL} $\to$ absent.
  \end{itemize}
  \muted{Cost $=$ number of nodes on the path $\le$ height $+1$. A short
    tree is a fast tree.}
\end{frame}

\begin{frame}[fragile]{Search in C}
  \scriptsize
\begin{verbatim}
struct node *bst_search(struct node *t, int key) {
    while (t != NULL && t->data != key) {
        if (key < t->data) t = t->left;
        else               t = t->right;
    }
    return t;          /* NULL means "not present" */
}
\end{verbatim}
  \vspace{0.1cm}
  \footnotesize
  \begin{itemize}
    \item Iterative, so no recursion stack: $O(1)$ extra space.
    \item Worst-case time is $O(h)$ where $h$ is the height. For a balanced
      tree $h=O(\log n)$; for a chain $h=n-1$.
    \item This is the same loop as binary search on a sorted array, with
      \texttt{left}/\texttt{right} pointers instead of index arithmetic.
  \end{itemize}
\end{frame}

\begin{frame}{Insert: Walk Down, Hang a Leaf}
  \begin{block}{Rule}
    Search for the key. If it is already present, reject it (or count it).
    If the search falls off the tree at \texttt{NULL}, put the new node
    there.
  \end{block}
  \begin{itemize}
    \item Insert $50,30,70,20,40,60,80$ into an empty tree:
      \begin{itemize}
        \item 50 becomes the root;
        \item 30 $<50$ $\to$ left of 50; 70 $>50$ $\to$ right of 50;
        \item 20 $<50,<30$ $\to$ left of 30; 40 $<50,>30$ $\to$ right of 30;
        \item 60 $>50,<70$ $\to$ left of 70; 80 $>50,>70$ $\to$ right of 70.
      \end{itemize}
    \item A new key always becomes a \key{leaf}. Insert never rearranges the
      existing nodes --- and that is exactly the weakness we will fix in
      Section 5.
  \end{itemize}
\end{frame}

\begin{frame}{Minimum and Maximum: Go All the Way Left or Right}
  \begin{itemize}
    \item \key{Minimum}: start at the root and keep taking \texttt{left}
      until there is no left child. In our tree: $50\to30\to20$, so the
      minimum is 20.
    \item \key{Maximum}: keep taking \texttt{right}: $50\to70\to80$, so the
      maximum is 80.
    \item Why it works: every key in a left subtree is smaller, so the
      smallest key cannot have anything to its left.
  \end{itemize}
  \begin{exampleblock}{Cost}
    Both are $O(h)$ --- one root-to-leaf walk. They are the first ordered
    operations a hash table cannot do at all.
  \end{exampleblock}
\end{frame}

\begin{frame}{Deletion, Cases 1 and 2: Zero or One Child}
  \begin{itemize}
    \item \key{Case 1 --- leaf}: remove it and set the parent's pointer to
      \texttt{NULL}. Delete 20 from our tree: 30's left pointer becomes
      \texttt{NULL}.
    \item \key{Case 2 --- one child}: splice the node out: the parent points
      straight to the node's only child. If 30 has only the right child 40,
      deleting 30 makes 40 the left child of 50.
    \item Both cases preserve the BST rule because the subtree that replaces
      the deleted node was already on the correct side.
  \end{itemize}
  \muted{Always \texttt{free} the removed node after saving its child
    pointer; otherwise the deletion leaks memory.}
\end{frame}

\begin{frame}{Deletion, Case 3: Two Children}
  \begin{center}
    % Before: delete 50. Successor = leftmost node of the right subtree = 60.
    % After: 60 replaces 50; the old 60 node is removed.
    \begin{tikzpicture}[every node/.style={font=\scriptsize}]
      \foreach \x/\y/\v in {1.2/2.2/50, 0.45/1.3/30, 1.95/1.3/70,
                            0.05/0.4/20, 0.85/0.4/40, 1.55/0.4/60, 2.35/0.4/80}
        \node[tnodesm] at (\x,\y) {\v};
      \draw[tedge] (1.2,2.2)--(0.45,1.3); \draw[tedge] (1.2,2.2)--(1.95,1.3);
      \draw[tedge] (0.45,1.3)--(0.05,0.4); \draw[tedge] (0.45,1.3)--(0.85,0.4);
      \draw[tedge] (1.95,1.3)--(1.55,0.4); \draw[tedge] (1.95,1.3)--(2.35,0.4);
      \node[tnodesmg] at (1.55,0.4) {60};
      \node[text=neutral] at (1.2,-0.35) {before: delete 50};

      \foreach \x/\y/\v in {6.0/2.2/60, 5.25/1.3/30, 6.75/1.3/70,
                            4.85/0.4/20, 5.65/0.4/40, 7.15/0.4/80}
        \node[tnodesm] at (\x,\y) {\v};
      \draw[tedge] (6.0,2.2)--(5.25,1.3); \draw[tedge] (6.0,2.2)--(6.75,1.3);
      \draw[tedge] (5.25,1.3)--(4.85,0.4); \draw[tedge] (5.25,1.3)--(5.65,0.4);
      \draw[tedge] (6.75,1.3)--(7.15,0.4);
      \node[tnodesmg] at (6.0,2.2) {60};
      \node[text=neutral] at (6.0,-0.35) {after: 60 replaces 50};
    \end{tikzpicture}
  \end{center}
  \vspace{-0.1cm}
  \footnotesize
  \begin{itemize}
    \item Find the \key{successor}: the smallest key in the right subtree
      (go right once, then left as far as possible). Here it is 60.
    \item Copy 60 into the node, then delete the old 60 node --- which has
      at most one child, so Case 1 or 2 applies.
  \end{itemize}
\end{frame}

\begin{frame}[fragile]{Deletion in C}
  \scriptsize
\begin{verbatim}
struct node *bst_delete(struct node *t, int key) {
    if (t == NULL) return NULL;
    if (key < t->data)      t->left  = bst_delete(t->left,  key);
    else if (key > t->data) t->right = bst_delete(t->right, key);
    else {                                  /* found it */
        if (t->left == NULL)  { struct node *r = t->right; free(t); return r; }
        if (t->right == NULL) { struct node *l = t->left;  free(t); return l; }
        struct node *s = t->right;           /* successor */
        while (s->left != NULL) s = s->left;
        t->data  = s->data;
        t->right = bst_delete(t->right, s->data);
    }
    return t;
}
\end{verbatim}
  \vspace{-0.1cm}
  \footnotesize
  \muted{The recursive call in the two-child case removes the successor,
    which has no left child. Every node is visited once: $O(h)$.}
\end{frame}

\begin{frame}{The Cost of a BST Is Its Height}
  \begin{center}
  {\footnotesize
  \begin{tabular}{l c c}
    \toprule
    \textbf{shape} & \textbf{height $h$} & \textbf{search / insert / delete} \\
    \midrule
    balanced (best) & $\lfloor \log_2 n \rfloor$ & $O(\log n)$ \\
    random insertions (average) & $\approx 1.39\log_2 n$ & $O(\log n)$ expected \\
    chain (worst) & $n-1$ & $O(n)$ \\
    \bottomrule
  \end{tabular}}
  \end{center}
  \begin{itemize}
    \item Every BST operation is one root-to-leaf walk, so all three cost
      $O(h)$.
    \item The average row assumes keys arrive in random order. Real data is
      often sorted --- names, timestamps, roll numbers --- and that is the
      dangerous case.
  \end{itemize}
\end{frame}

\begin{frame}{When Inserts Are Sorted, the Tree Becomes a List}
  \begin{center}
    % Left: inserting 10,20,30,40,50 in order makes a right chain (height 4).
    % Right: a balanced tree over the same keys has height 2.
    \begin{tikzpicture}[every node/.style={font=\scriptsize}]
      \foreach \y/\v in {3.2/10, 2.4/20, 1.6/30, 0.8/40, 0/50}
        \node[tnodesm] at (0,\y) {\v};
      \foreach \y in {2.8,2.0,1.2,0.4}
        \draw[tedge] (0,\y+0.4) -- (0,\y-0.4);
      \node[text=negative, align=center] at (0,-0.65)
        {\bad{height 4}\\ search(50) $=5$ comparisons};

      \foreach \x/\y/\v in {4.5/2.4/30, 3.8/1.4/20, 5.2/1.4/40,
                            3.45/0.4/10, 4.15/0.4/25, 4.85/0.4/35, 5.55/0.4/50}
        \node[tnodesm] at (\x,\y) {\v};
      \draw[tedge] (4.5,2.4)--(3.8,1.4); \draw[tedge] (4.5,2.4)--(5.2,1.4);
      \draw[tedge] (3.8,1.4)--(3.45,0.4); \draw[tedge] (3.8,1.4)--(4.15,0.4);
      \draw[tedge] (5.2,1.4)--(4.85,0.4); \draw[tedge] (5.2,1.4)--(5.55,0.4);
      \node[text=positive, align=center] at (4.5,-0.65)
        {\good{height 2}\\ search(50) $=3$ comparisons};
    \end{tikzpicture}
  \end{center}
  \vspace{-0.2cm}
  \begin{itemize}
    \item Insert $10,20,30,40,50$ in order: every new key is larger, so it
      always goes right. The ``tree'' is a linked list with extra steps.
    \item Nothing in the BST insert rule prevents this. We need a rule that
      repairs the shape.
  \end{itemize}
\end{frame}

\transitionslide{A tree that repairs itself.}

% =============================================================================
% ACT 5 --- AVL TREES
% =============================================================================
\section{AVL Trees}

\sectiondivider{Section 5}{Balance Factors and Rotations}

\begin{frame}{The Balance Rule}
  \begin{block}{Definition --- AVL tree}
    An \key{AVL tree} is a BST in which, for every node, the heights of the
    left and right subtrees differ by at most $1$.
  \end{block}
  \begin{itemize}
    \item Define the \key{balance factor} of a node:
      \[
        \mathrm{bf}(x) = \mathrm{height}(\text{left child of }x)
                       - \mathrm{height}(\text{right child of }x).
      \]
    \item The AVL rule: $\mathrm{bf}(x) \in \{-1,0,+1\}$ for every node.
    \item If an insert or delete makes some $|\mathrm{bf}| = 2$, we repair
      the tree with a \key{rotation}. A rotation is a local pointer change;
      it keeps the inorder order.
  \end{itemize}
  \muted{Named after Adelson-Velsky and Landis (1962). Height convention:
    leaf $=0$, empty $=-1$. Definitions and rotations:
    \cite[pp.~142--149]{Wirth2004_algorithms_data_structures}.}
\end{frame}

\begin{frame}{Balance Factors on a Picture}
  \begin{center}
    \begin{tikzpicture}[every node/.style={font=\small}]
      \node[tnode] (n50) at (0,2.5) {50};
      \node[tnode] (n30) at (-1.3,1.4) {30};
      \node[tnode] (n70) at (1.3,1.4) {70};
      \node[tnode] (n20) at (-1.9,0.3) {20};
      \draw[tedge] (n50)--(n30); \draw[tedge] (n50)--(n70);
      \draw[tedge] (n30)--(n20);
      \node[text=primary-gold, font=\scriptsize] at (0.0,3.38) {bf $=+1$};
      \node[text=primary-gold, font=\scriptsize] at (-1.3,2.30) {bf $=+1$};
      \node[text=primary-gold, font=\scriptsize] at (1.3,2.30) {bf $=0$};
      \node[text=primary-gold, font=\scriptsize] at (-1.9,-0.45) {bf $=0$};
    \end{tikzpicture}
  \end{center}
  \vspace{-0.1cm}
  \begin{itemize}
    \item Heights: 20 and 70 are leaves (height 0); 30 has height 1;
      50 has height 2.
    \item $\mathrm{bf}(50)=+1$: its left subtree (30, with 20) has height
      1; its right subtree (70) has height 0.
    \item $\mathrm{bf}(30)=+1$: 20 has height 0, the missing right child
      has height $-1$. And $\mathrm{bf}(20)=\mathrm{bf}(70)=0$.
    \item Every factor is in $\{-1,0,+1\}$: this is a valid AVL tree.
  \end{itemize}
\end{frame}

\begin{frame}{A Rotation Is Three Pointer Moves}
  \begin{center}
    % Before: y=30 at (1,2.2), x=20 at (0,1.2), T1 below x, T2 right of x.
    % After: x=20 becomes root; y=30 becomes its right child; T2 moves to y.left.
    \begin{tikzpicture}[every node/.style={font=\scriptsize}]
      \node[tnodesm] (y) at (1,2.2) {30};
      \node[tnodesm] (x) at (0,1.2) {20};
      \node[tnodesm] (t1) at (-0.5,0.2) {10};
      \node[draw=neutral, fill=white, minimum width=0.55cm,
        minimum height=0.45cm] (t2) at (0.5,0.2) {T2};
      \draw[tedge] (y)--(x); \draw[tedge] (x)--(t1); \draw[tedge] (x)--(t2);
      \node[text=neutral] at (0.5,-0.55) {before};

      \draw[-{Latex[length=2.5mm]}, primary-gold, thick]
        (2.0,1.2) -- (2.9,1.2);

      \node[tnodesm] (X) at (5.0,2.2) {20};
      \node[tnodesm] (T1) at (4.35,1.2) {10};
      \node[tnodesm] (Y) at (5.65,1.2) {30};
      \node[draw=neutral, fill=white, minimum width=0.55cm,
        minimum height=0.45cm] (T2b) at (5.0,0.2) {T2};
      \draw[tedge] (X)--(T1); \draw[tedge] (X)--(Y); \draw[tedge] (Y)--(T2b);
      \node[text=neutral] at (5.0,-0.55) {after right rotation at 30};
    \end{tikzpicture}
  \end{center}
  \vspace{-0.1cm}
  \begin{itemize}
    \item The middle key (20) rises; the old root (30) becomes its right
      child; 20's old right subtree (T2) becomes 30's left subtree.
    \item Inorder before and after: $10, 20, \mathrm{T2}, 30$ --- unchanged.
      A rotation rearranges the shape, never the order.
  \end{itemize}
\end{frame}

\begin{frame}{The Four Cases}
  \begin{center}
  {\footnotesize
  \begin{tabular}{l l l}
    \toprule
    \textbf{case} & \textbf{path from the unbalanced node} & \textbf{repair} \\
    \midrule
    LL & left child, then left child  & right rotation once \\
    RR & right child, then right child & left rotation once \\
    LR & left child, then right child & left rotate child, then right rotate node \\
    RL & right child, then left child & right rotate child, then left rotate node \\
    \bottomrule
  \end{tabular}}
  \end{center}
  \begin{itemize}
    \item Read the name as the \emph{path}, not the first rotation. LR means
      ``went Left, then Right.''
    \item The first rotation turns LR into LL; the second rotation finishes
      the job. That is why LR and RL need two rotations.
    \item Only the nodes on the path from the inserted leaf back to the root
      can change balance. Fix the lowest violation first.
  \end{itemize}
\end{frame}

\begin{frame}{LL and RR: One Rotation Each}
  \begin{center}
    \begin{tikzpicture}[every node/.style={font=\scriptsize}]
      % LL: 30-20-10 -> right rotation at 30 -> 20 root.
      \node[tnodesm] (l1) at (0,2.2) {30};
      \node[tnodesm] (l2) at (0,1.2) {20};
      \node[tnodesm] (l3) at (0,0.2) {10};
      \draw[tedge] (l1)--(l2); \draw[tedge] (l2)--(l3);
      \node[text=neutral, font=\tiny] at (-0.15,-0.45) {LL: 30--20--10};
      \draw[-{Latex[length=2mm]}, primary-gold, thick] (0.55,1.2)--(1.15,1.2);
      \node[tnodesm] (L1) at (1.8,2.2) {20};
      \node[tnodesm] (L2) at (1.3,1.2) {10};
      \node[tnodesm] (L3) at (2.3,1.2) {30};
      \draw[tedge] (L1)--(L2); \draw[tedge] (L1)--(L3);
      \node[text=positive, font=\tiny] at (1.95,-0.45) {\good{right rotate}};

      % RR: 10-20-30 -> left rotation at 10 -> 20 root.
      \node[tnodesm] (r1) at (4.6,2.2) {10};
      \node[tnodesm] (r2) at (4.6,1.2) {20};
      \node[tnodesm] (r3) at (4.6,0.2) {30};
      \draw[tedge] (r1)--(r2); \draw[tedge] (r2)--(r3);
      \node[text=neutral, font=\tiny] at (4.45,-0.45) {RR: 10--20--30};
      \draw[-{Latex[length=2mm]}, primary-gold, thick] (5.15,1.2)--(5.75,1.2);
      \node[tnodesm] (R1) at (6.4,2.2) {20};
      \node[tnodesm] (R2) at (5.9,1.2) {10};
      \node[tnodesm] (R3) at (6.9,1.2) {30};
      \draw[tedge] (R1)--(R2); \draw[tedge] (R1)--(R3);
      \node[text=positive, font=\tiny] at (6.55,-0.45) {\good{left rotate}};
    \end{tikzpicture}
  \end{center}
  \vspace{-0.2cm}
  \begin{itemize}
    \item LL and RR are mirror images. One rotation lifts the middle key to
      the top.
    \item After the rotation the tree is a perfect three-node tree of height
      1, so the balance factor at the old root is back in range.
  \end{itemize}
\end{frame}

\begin{frame}{LR: Left Rotate the Child, Then Right Rotate the Node}
  \begin{center}
    \begin{tikzpicture}[every node/.style={font=\scriptsize}]
      % Step 0: 30 -> left 10 -> right 20.
      \node[tnodesm] (a1) at (0,2.2) {30};
      \node[tnodesm] (a2) at (-0.7,1.2) {10};
      \node[tnodesm] (a3) at (0.0,0.2) {20};
      \draw[tedge] (a1)--(a2); \draw[tedge] (a2)--(a3);
      \node[text=neutral, font=\tiny] at (0,-0.45) {start};

      \draw[-{Latex[length=2mm]}, primary-gold, thick] (0.6,1.2)--(1.2,1.2);
      % Step 1: left rotate 10 -> 30 -> left 20 -> left 10.
      \node[tnodesm] (b1) at (2.1,2.2) {30};
      \node[tnodesm] (b2) at (1.5,1.2) {20};
      \node[tnodesm] (b3) at (1.0,0.2) {10};
      \draw[tedge] (b1)--(b2); \draw[tedge] (b2)--(b3);
      \node[text=neutral, font=\tiny] at (2.1,-0.45) {left rotate 10};

      \draw[-{Latex[length=2mm]}, primary-gold, thick] (2.7,1.2)--(3.3,1.2);
      % Step 2: right rotate 30 -> 20 root.
      \node[tnodesm] (c1) at (4.2,2.2) {20};
      \node[tnodesm] (c2) at (3.7,1.2) {10};
      \node[tnodesm] (c3) at (4.7,1.2) {30};
      \draw[tedge] (c1)--(c2); \draw[tedge] (c1)--(c3);
      \node[text=positive, font=\tiny] at (4.2,-0.45) {\good{right rotate 30}};
    \end{tikzpicture}
  \end{center}
  \vspace{-0.2cm}
  \begin{itemize}
    \item The first rotation makes the path LL. The second rotation is the
      LL repair we already know.
    \item Inorder stays $10,20,30$ at every step.
  \end{itemize}
\end{frame}

\begin{frame}{RL: Right Rotate the Child, Then Left Rotate the Node}
  \begin{center}
    \begin{tikzpicture}[every node/.style={font=\scriptsize}]
      % Step 0: 10 -> right 30 -> left 20.
      \node[tnodesm] (a1) at (0,2.2) {10};
      \node[tnodesm] (a2) at (0.7,1.2) {30};
      \node[tnodesm] (a3) at (0.0,0.2) {20};
      \draw[tedge] (a1)--(a2); \draw[tedge] (a2)--(a3);
      \node[text=neutral, font=\tiny] at (0,-0.45) {start};

      \draw[-{Latex[length=2mm]}, primary-gold, thick] (0.6,1.2)--(1.2,1.2);
      % Step 1: right rotate 30 -> 10 -> right 20 -> right 30.
      \node[tnodesm] (b1) at (2.1,2.2) {10};
      \node[tnodesm] (b2) at (2.7,1.2) {20};
      \node[tnodesm] (b3) at (3.2,0.2) {30};
      \draw[tedge] (b1)--(b2); \draw[tedge] (b2)--(b3);
      \node[text=neutral, font=\tiny] at (2.1,-0.45) {right rotate 30};

      \draw[-{Latex[length=2mm]}, primary-gold, thick] (2.7,1.2)--(3.3,1.2);
      % Step 2: left rotate 10 -> 20 root.
      \node[tnodesm] (c1) at (4.2,2.2) {20};
      \node[tnodesm] (c2) at (3.7,1.2) {10};
      \node[tnodesm] (c3) at (4.7,1.2) {30};
      \draw[tedge] (c1)--(c2); \draw[tedge] (c1)--(c3);
      \node[text=positive, font=\tiny] at (4.2,-0.45) {\good{left rotate 10}};
    \end{tikzpicture}
  \end{center}
  \vspace{-0.2cm}
  \begin{itemize}
    \item Mirror of LR: first turn RL into RR, then do the single left
      rotation.
    \item Both double rotations restore the same three-node shape. The
      difference is only which key was inserted last.
  \end{itemize}
\end{frame}

\begin{frame}{Watch an Insert Rebalance: 30, 20, 10, 25, 28}
  \begin{columns}[T]
    \begin{column}{0.55\textwidth}
      \footnotesize
      \begin{enumerate}
        \item Insert 30, 20, 10: LL at 30. Right-rotate 30; \textbf{20}
          becomes root, with 10 left and 30 right.
        \item Insert 25: it goes left of 30. Factors still valid:
          $\mathrm{bf}(30)=+1$, $\mathrm{bf}(20)=-1$.
        \item Insert 28: it goes right of 25. Now $\mathrm{bf}(30)=+2$;
          the path from 30 is left (25) then right (28), so this is
          \textbf{LR}.
        \item Left-rotate 25, then right-rotate 30. \textbf{28} becomes
          the root of that subtree.
      \end{enumerate}
    \end{column}
    \begin{column}{0.42\textwidth}
      \centering
      \begin{tikzpicture}[every node/.style={font=\scriptsize}]
        \node[tnodesm] (r) at (0,2.4) {20};
        \node[tnodesm] (l) at (-1.0,1.3) {10};
        \node[tnodesm] (m) at (1.0,1.3) {28};
        \node[tnodesm] (ml) at (0.5,0.2) {25};
        \node[tnodesm] (mr) at (1.5,0.2) {30};
        \draw[tedge] (r)--(l); \draw[tedge] (r)--(m);
        \draw[tedge] (m)--(ml); \draw[tedge] (m)--(mr);
      \end{tikzpicture}
      \\[-0.1cm]
      {\scriptsize final tree: height 2}
    \end{column}
  \end{columns}
\end{frame}

\begin{frame}[fragile]{AVL Insert in C: Rotate, Then Recompute Heights}
  \scriptsize
\begin{verbatim}
int height(struct node *t) { return t ? t->height : -1; }
int bf(struct node *t)     { return t ? height(t->left) - height(t->right) : 0; }

struct node *rotate_right(struct node *y) {
    struct node *x  = y->left;          /* x rises */
    struct node *t2 = x->right;         /* T2 changes parent */
    x->right = y;  y->left = t2;
    y->height = 1 + max(height(y->left),  height(y->right));
    x->height = 1 + max(height(x->left),  height(x->right));
    return x;                            /* new subtree root */
}
/* rotate_left is the mirror image. */
\end{verbatim}
  \vspace{-0.1cm}
  \footnotesize
  \begin{itemize}
    \item After a normal BST insert, walk back up. At each node update
      \texttt{height}, compute \texttt{bf}, and apply the LL/RR/LR/RL case
      whose factor is $\pm2$.
    \item Store the height in the node so the check is $O(1)$ per node,
      not another tree walk.
  \end{itemize}
\end{frame}

\begin{frame}{The Height Guarantee: Why AVL Is $O(\log n)$}
  \begin{itemize}
    \item Let $N(h)$ be the \emph{fewest} nodes an AVL tree of height $h$
      can have. The worst shape has one subtree of height $h-1$ and the
      other of height $h-2$:
      \[
        N(h) = N(h-1) + N(h-2) + 1,\qquad N(0)=1,\ N(1)=2.
      \]
    \item This is the Fibonacci recurrence (shifted), so $N(h)$ grows
      exponentially: $1, 2, 4, 7, 12, 20, 33, \dots$
    \item Therefore $h = O(\log n)$, with the precise bound
      $h \le 1.44\log_2(n+2) - 0.33$.
    \item In practice: $n=1{,}000{,}000$ gives $h\le 28$, against
      $999{,}999$ for the chain. Search, insert and delete are all
      $O(\log n)$ worst case.
  \end{itemize}
  \muted{You do not need the constant $1.44$ in an exam; you do need the
    recurrence and the conclusion.}
\end{frame}

\begin{frame}{Deletion in an AVL Tree}
  \begin{itemize}
    \item Delete exactly as in a BST (the three cases), then walk from the
      parent of the removed node back up to the root.
    \item At each node, recompute the height and balance factor. If
      $|\mathrm{bf}|=2$, apply the same four cases.
    \item Deletion can need \emph{several} rotations, not just one: removing
      a node can unbalance several ancestors. Insert needs at most one
      repair on the path.
    \item Still $O(\log n)$: the walk has $\log n$ levels and each rotation
      is $O(1)$.
  \end{itemize}
  \begin{exampleblock}{The invariant to remember}
    A rotation preserves the BST rule and the inorder sequence. That is the
    only reason we are allowed to repair the shape without touching the data.
  \end{exampleblock}
\end{frame}

\begin{frame}{Choosing: Tree, Hash Table, or Sorted Array?}
  \begin{center}
  {\footnotesize
  \begin{tabular}{l c c c c}
    \toprule
    \textbf{structure} & \textbf{search} & \textbf{insert} & \textbf{ordered ops} & \textbf{space} \\
    \midrule
    sorted array & $O(\log n)$ & $O(n)$ & yes & $O(n)$ \\
    hash table   & $O(1)$ expected & $O(1)$ expected & no & $O(n)$ \\
    plain BST    & $O(h)$ & $O(h)$ & yes & $O(n)$ \\
    AVL tree     & $O(\log n)$ & $O(\log n)$ & yes & $O(n)$ \\
    \bottomrule
  \end{tabular}}
  \end{center}
  \begin{itemize}
    \item Point lookup only, order irrelevant $\to$ hash table.
    \item Sorted output, next/previous, range queries, min/max $\to$ AVL.
    \item Static data read many times and never updated $\to$ sorted array.
  \end{itemize}
  \muted{``Trees are better than hashing'' is false; ``trees are better when
    order is part of the question'' is true.}
\end{frame}

% =============================================================================
% ACT 6 --- CLOSE + UNIT III REVIEW
% =============================================================================
\section{Review and Close}

\begin{frame}{Three Classic Mistakes}
  \begin{itemize}
    \item \bad{``It is a binary tree, so search is $O(\log n)$.''} Only a
      \emph{short} BST is logarithmic. Inserting sorted keys builds a chain
      and search becomes $O(n)$.
    \item \bad{``Inorder is always sorted.''} Only for a BST. On an arbitrary
      binary tree, inorder is just left--node--right.
    \item \bad{``A rotation changes the keys.''} A rotation changes pointers
      and heights, never the inorder sequence. If the inorder changes, the
      rotation is wrong.
  \end{itemize}
  \begin{exampleblock}{How to revise}
    For each structure, write four lines: search cost, insert cost, whether
    order survives, and the one case that breaks it. Then trace one insert
    and one delete on a seven-node tree, and one LL or RR rotation.
  \end{exampleblock}
\end{frame}

\begin{frame}{Unit III in One Page}
  \begin{center}
  {\scriptsize
  \begin{tabular}{c l l l}
    \toprule
    \textbf{week} & \textbf{idea} & \textbf{key cost} & \textbf{key trap} \\
    \midrule
    8 & search; elementary sorts &
      $\log n$ search; $n^2$ sorts & binary search needs sorted data \\
    9 & efficient sorts &
      $n\log n$; radix $d(n+b)$ & quicksort is $n^2$ on sorted input \\
    10 & hashing &
      $O(1)$ expected lookup & order is lost; tombstone on delete \\
    11 & trees, BST, AVL &
      $O(\log n)$ if balanced & a plain BST can become a chain \\
    \bottomrule
  \end{tabular}}
  \end{center}
  \begin{itemize}
    \item The unit's single theme: \key{the cost of an operation is decided
      by the shape of the data, not by the name of the structure}.
    \item Week 12 removes the last restriction: trees have one root and no
      cycles; graphs may have many paths, cycles, and no root at all.
  \end{itemize}
\end{frame}

\begin{frame}{Week 11 Summary}
  \begin{itemize}
    \item \key{Tree vocabulary} --- root, leaf, parent, child, sibling,
      subtree; depth counts down from the root, height counts up from a
      leaf. Height controls the cost.
    \item \key{Traversals} --- inorder (left, node, right), preorder (node,
      left, right), postorder (left, right, node), and level order (queue).
      Each is $O(n)$; each has a job.
    \item \key{Expression trees} --- operands are leaves, operators are
      internal nodes; postorder evaluates, preorder gives prefix, inorder
      gives infix (with parentheses).
    \item \key{BST} --- left subtree $<$ node $<$ right subtree at every
      node. Search, insert, min, max and delete are $O(h)$; delete has
      three cases, and the two-child case uses the successor.
    \item \key{AVL} --- keep every balance factor in $\{-1,0,+1\}$; repair
      with LL/RR (one rotation) or LR/RL (two). Height, and therefore every
      operation, is $O(\log n)$.
  \end{itemize}
  \begin{exampleblock}{Next week}
    Graphs: many connections, cycles, and no single root. BFS and DFS,
    Prim and Kruskal, and Dijkstra's shortest paths.
  \end{exampleblock}
\end{frame}

\begin{frame}{Before You Go: Predict}
  \begin{block}{Three one-line predictions}
    Answer from the rules, then check in the lab.
  \end{block}
  \begin{enumerate}
    \item Insert $5,3,8,1,4,7,9$ into an empty BST. What is the inorder
      output, and what is the height?
    \item Delete 5 from that tree. Which key replaces it, and why?
    \item Insert $10,20,30$ into an empty AVL tree. Draw the result and name
      the rotation.
  \end{enumerate}
  \muted{Answers: inorder $1,3,4,5,7,8,9$; height 2. Deleting 5 uses its
    successor 7 (smallest key in the right subtree). The AVL insert is the
    RR case: one left rotation at 10 gives 20 as root, with 10 left and 30
    right.}
\end{frame}

\begin{frame}[allowframebreaks]{References}
  \tiny
  \bibliographystyle{plain}
  \bibliography{../../Bibliography_base}
\end{frame}

\end{document}
